\chapter{现代光学基础}

\problem{8.1}{（1）计算氢原子最低的四个能级的能量大小，并把它们画成能级图；（2）计算这四个能级之间跃迁的最小的频率是多少？}
\answerarea

\problem{8.2}{当玻尔描述的氢原子从 $n=2$ 的轨道跃迁到 $n=1$ 的轨道后，问（1）轨道半径变化了多少？（2）能量改变了多少？}
\answerarea
\clearpage

\problem{8.3}{一个电子受激跃迁到氢原子的第三玻尔轨道，试计算回到基态时有可能发射的光子的波长。}
\answerarea

\problem{8.4}{氦氖激光的单色性为 $6\times 10^{-10}$，则其相干长度为多少？}
\answerarea
\clearpage

\problem{8.5}{证明当每个模内的平均光子数大于 $1$ 时，腔内振荡以受激辐射为主。（提示：以受激辐射为主时，$\dfrac{w_{21}}{A_{21}}=\dfrac{B_{21}\rho}{A_{21}}>1$。）}
\answerarea

\problem{8.6}{一质地均匀的材料对光的吸收系数为 $0.01\,\mathrm{mm^{-1}}$，光通过 $10\,\mathrm{cm}$ 长的该材料后，光强损耗多少？若对 $10\,\mathrm{cm}$ 长的另一种材料，出射光强是入射光强的 $2.718$ 倍，问该材料的增益系数为多少？}
\answerarea
\clearpage

\problem{8.7}{如激光器分别在波长为 $10\,\mathrm{\mu m}$ 和 $500\,\mathrm{nm}$ 时输出 $1\,\mathrm{W}$ 的连续功率，求每秒从激光上能级向下能级跃迁时对应的粒子数。}
\answerarea

\problem{8.8}{设一对激光能级为 $E_2$ 和 $E_1$，上下能级上的粒子密度分别为 $n_2$ 和 $n_1$，两能级间的跃迁频率为 $\nu$，波长为 $\lambda$。试求（1）当 $\nu=3\,000\,\mathrm{MHz}$ 和 $\lambda=1\,\mathrm{\mu m}$ 时，$n_2/n_1$ 分别为多少？假定此时 $T=300\,\mathrm{K}$；（2）当 $\lambda=1\,\mathrm{\mu m}$，$n_2/n_1=0.1$ 时，温度 $T$ 为多少？}
\answerarea
\clearpage

\problem{8.9}{假若受激辐射爱因斯坦系数 $B=10^{19}\,\mathrm{m^3/(W\cdot s^3)}$，试对下列波长计算自发辐射跃迁系数和自发辐射寿命。（1）$\lambda=6\,\mathrm{\mu m}$；（2）$\lambda=600\,\mathrm{nm}$；（3）$\lambda=60\,\mathrm{nm}$；（4）$\lambda=0.6\,\mathrm{nm}$。如果光强为 $10\,\mathrm{W/mm^2}$，求受激跃迁速率。}
\answerarea

\problem{8.10}{有一个二能级系统，如果能级差为 $1.602\times 10^{-21}\,\mathrm{J}$，求：（1）当 $T=10^2\,\mathrm{K}$、$10^5\,\mathrm{K}$ 和 $10^8\,\mathrm{K}$ 时，上下两能级粒子数 $N_2$ 与 $N_1$ 之比；（2）如 $N_2=N_1$，则相当于多高的温度 $\mathrm{K}$；（3）要有怎样的温度，粒子数才能发生反转；（4）如果用“负温度”的概念来描述粒子数反转的状态，则 $T=-10^4\,\mathrm{K}$ 和 $T=-10^8\,\mathrm{K}$ 两个温度中哪一个高？$T=-10^5\,\mathrm{K}$ 和 $T=+10^8\,\mathrm{K}$ 中哪一个高？}
\answerarea
\clearpage

\problem{8.11}{衍射和傅里叶变换的关系是什么？如何应用傅里叶变换来计算菲涅耳衍射和夫琅禾费衍射？给出计算的方法和步骤。}
\answerarea

\problem{8.12}{从惠更斯-菲涅耳原理的数学表达式出发，
\begin{enumerate}
  \item 求出菲涅耳衍射方程；
  \item 证明菲涅耳衍射的作用相当于空间不变的线性系统；
  \item 并求出传递函数；
  \item 求出夫琅禾费衍射方程。
\end{enumerate}}
\answerarea
\clearpage

\problem{8.13}{描述二元振幅滤波器的种类及其应用。}
\answerarea

\problem{8.14}{说明泽尼克相衬法中采用的滤波系统，并计算像面上的强度分布情况。}
\answerarea
\clearpage

\problem{8.15}{题 8.15 图为电寻址液晶光阀试验系统，主要由激光光源、准直扩束系统、高分辨率电寻址透射式液晶光阀、傅里叶变换透镜、CCD 探测器等构成。与光寻址液晶光阀不同的是，该系统可以直接把计算机产生的信息显示在液晶光阀上，可以实时地进行图像处理，且方便实验操作。

请解释该系统的工作原理。\figcap[0.74]{8_15}{题 8.15 图\quad 电寻址液晶光阀实验系统}}
\answerarea

\sproblem{8.16}{研究性课题：通过书面材料或互联网，查找激光技术的进展，或激光的应用等方面的知识，写出报告，在同学之间相互交流。}
\answerarea
\clearpage
